\documentclass[10pt]{article}
\PassOptionsToPackage{landscape}{geometry}
\usepackage{../common/d3report}
\hypersetup{pdftitle={CS455 D3 Requirements Traceability Matrix},pdfauthor={Mainak Sarkar}}
\newcommand{\TraceRequirements}{193}
\newcommand{\TraceJira}{79}
\newcommand{\TracePulls}{50}
\newcommand{\TraceCases}{654}
\newcommand{\TraceFiles}{115}
\newcommand{\TraceByReading}{53}
\newcommand{\TracePassing}{193}
\newcommand{\UnitRunDate}{2026-11-04}
\newcommand{\UnitRunCases}{1535}
\newcommand{\IntegrationRunDate}{2026-11-02}
\newcommand{\IntegrationRunCases}{109}
\newcommand{\SnapshotPulls}{87}
\newcommand{\CoverageStatements}{87.6}


% A shaded row that starts each area of the matrix.
\newcommand{\sectionrow}[1]{\rule{0pt}{2.6ex}{\sffamily\bfseries\color{hdrblue}#1}}

\begin{document}

\makereporttitle{Requirements Traceability Matrix}
  {Deliverable 3, Campus Ride-Pooling and Split-Fare System}
  {Prepared by Mainak Sarkar \quad Requirement $\rightarrow$ Jira $\rightarrow$ GitHub $\rightarrow$ test
  \quad Tests run on \UnitRunDate}

\keybox{\textbf{Result.} \TraceRequirements{} requirements of the Deliverable 1 specification are
traced to \TraceJira{} Jira issues, \TracePulls{} merged pull requests and \TraceCases{} automated
test cases in \TraceFiles{} test files. Every one of those test cases passed in the recorded runs:
the unit suite (\UnitRunCases{} tests) and the integration and adversarial suite
(\IntegrationRunCases{} tests against a real PostgreSQL database), both on the master branch.}

\section{Reading the matrix}
\begin{tabularx}{\linewidth}{@{}>{\bfseries}p{0.14\linewidth}X@{}}
\toprule
Requirement & The requirement's identifier and statement from the Deliverable 1 specification (rider,
ride owner and Operations Admin documents). Long statements are shortened; the full text of each is
in \file{traceability-matrix.csv}. \\
Jira & The CS455 Jira issues of the pull requests in the GitHub column, as named in their titles,
branches, own commits and Jira sections. \\
GitHub & The merged pull requests that introduced the code under test, that cite the requirement,
and that introduced its tests; up to five are shown, the CSV lists all of them. \\
Tests & The test files that verify the requirement, each with the number of its test cases that do,
largest first. A test verifies a requirement when its name cites the requirement's identifier.
Requirements marked \textsuperscript{\dag} are verified by tests whose names do not cite the
identifier; each of those links was made by reading the test and is listed in
\file{data/manual-links.json}. \\
Result & Passed test cases out of those linked, from the recorded runs. \\
\bottomrule
\end{tabularx}

\section{How the matrix is built}
Every cell comes from a recorded source, so the matrix can be rebuilt and checked rather than
maintained by hand. A requirement is listed when at least one passing automated test verifies it.

\begin{itemize}
  \item \textbf{Requirements}: parsed from the Deliverable 1 requirement documents in
    \texttt{Deliverables/D1/Requirements Engineering/users/}.
  \item \textbf{Tests}: the recorded unit run (\file{data/unit-cases.json}) and integration run
    (\file{Deliverables/D3/test-run/integration-results.json}). Test names such as
    ``complaints (FR-RD-13.1 to 13.4)'' or ``(OA-FR-34, 36)'' cite every requirement in the range or
    list. The build refuses to run if a link in \file{data/manual-links.json} matches no recorded test.
  \item \textbf{GitHub}: a snapshot of the \SnapshotPulls{} merged pull requests with their changed
    files and commits (\file{data/prs.json}). The pull request that introduced a file is the first
    merged one that changed it. A commit counts only for the first merged pull request that
    contains it, since a stacked pull request also lists the commits below it.
  \item \textbf{Coverage}: NFR-RD-17 is verified by the unit suite's statement coverage
    (\CoverageStatements\,\%) against the 80\,\% gate that continuous integration enforces.
\end{itemize}

\keybox{\small
\texttt{node Deliverables/D3/Traceability\_Matrix/snapshot-prs.mjs}\hfill (merged pull requests)\\
\texttt{node Deliverables/D3/Traceability\_Matrix/make-matrix.mjs -{}-unit <run.json> -{}-coverage <summary.json>}\hfill (tables and CSV)\\[2pt]
The commands that record the test runs are in \file{Deliverables/D3/Traceability_Matrix/README.md}.}

\section{Summary by area}
\begin{center}\small
\begin{tabular}{@{}lrrrr@{}}
\toprule
Area & Requirements & Jira issues & Pull requests & Test cases \\
\midrule
\tabinput{generated/section-summary.tex}
\midrule
\textbf{Total (distinct)} & \textbf{\TraceRequirements} & \textbf{\TraceJira} & \textbf{\TracePulls} & \textbf{\TraceCases} \\
\bottomrule
\end{tabular}
\end{center}

\section{The matrix}
{\scriptsize
\setlength{\tabcolsep}{4pt}
\renewcommand{\arraystretch}{1.15}
\begin{longtable}{@{}>{\raggedright\arraybackslash}p{0.39\linewidth}>{\raggedright\arraybackslash}p{0.14\linewidth}>{\raggedright\arraybackslash}p{0.12\linewidth}>{\raggedright\arraybackslash}p{0.24\linewidth}>{\raggedright\arraybackslash}p{0.065\linewidth}@{}}
\toprule
\textbf{Requirement} & \textbf{Jira} & \textbf{GitHub} & \textbf{Tests} & \textbf{Result} \\
\midrule
\endfirsthead
\toprule
\textbf{Requirement} & \textbf{Jira} & \textbf{GitHub} & \textbf{Tests} & \textbf{Result} \\
\midrule
\endhead
\midrule
\multicolumn{5}{@{}r@{}}{\textsuperscript{\dag}\,verified by tests whose names do not cite the identifier (see \file{data/manual-links.json})} \\
\endfoot
\bottomrule
\endlastfoot
\tabinput{generated/matrix-rows.tex}
\end{longtable}
}

\end{document}
